
\begin{frame}{So Which Definition Do You Use?}

    This is a decision about \textbf{harms}, not a hyperparameter search. Four questions, in order:

    \vspace{0.8em}

    \begin{enumerate}
        \setlength\itemsep{0.55em}
        \item \textbf{Who is affected, and what is the harm?} Being wrongly \emph{denied}
              something (allocative) or being wrongly \emph{flagged} (punitive) point to different
              error rates.
        \item \textbf{Do you trust the label $Y$?} If the label encodes past discrimination,
              definitions that condition on $Y$ inherit it --- lean towards independence.
        \item \textbf{Is the score read by a human?} If yes, calibration within groups is close to
              mandatory, or you are misinforming the decision-maker.
        \item \textbf{Can you measure it?} A criterion you cannot estimate --- no group labels, or
              slices too small --- is a wish, not a criterion (see the next section).
    \end{enumerate}

    \vspace{0.9em}
    \begin{block}{}
        \centering
        Write the chosen criterion, and the criteria you consciously gave up, into the model card.
        An unstated choice is still a choice.
    \end{block}
\end{frame}
